\chapter{Родительская масса бозонного дефекта и проверка портала Хиггса}

Предыдущая проверка свела локализованный материальный остаток к нейтральному
скалярному бозону без внутреннего ротора. Теперь необходимо отделить три
разных утверждения: допустимость лоренц-ковариантного действия, вычисление
безразмерного профиля и вывод абсолютных массы, радиуса и наблюдаемой связи.

\section{Проектная и литературная граница}

Для моделей Фаддеева--Скирма коэффициенты двух- и четырёхпроизводных членов
задают единицы длины и энергии; две комбинации коэффициентов можно удалить
только выбором этих единиц, но не получить из одной топологии. Это явно
показано в \path{arXiv:1311.4454}. В выводах из высшеразмерной теории
Янга--Миллса оба члена действительно возникают совместно, однако их
коэффициенты являются интегралами по дополнительному профилю; см.
\path{arXiv:1807.11803}. Следовательно, наличие правильной составной
связности ещё не заменяет родительский вывод масштаба.

Калибровочно-синглетное скалярное поле может иметь перенормируемую связь с
$H^\dagger H$; это стандартный портал Хиггса, см.
\path{arXiv:hep-ph/0605188}. Но симметрийная допустимость оператора не
фиксирует его коэффициент и тем более не доказывает, что он присутствует в
конкретном родительском функционале.

Ранее Том V уже установил независимость спектральных моментов $f_2,f_0$,
масштаба отсечения и внутренней длины. Настоящая проверка не повторяет этот
запрет абстрактно, а применяет его к явному конечному профилю Тома VI и к
единому функционалу $M_{300}$.

\section{Допустимое лоренцево продолжение}

Статические инварианты допускают естественную запись в пространстве-времени:
\begin{equation}
 \mathcal L_Q=
 \frac{Z_Q}{2}\Tr(D_\mu QD^\mu Q)
 -\frac{Z_F}{4}\Tr(F_{\mu\nu}F^{\mu\nu})
 -V(Q).
 \label{eq:v6-bosonic-lorentz-admissible-action}
\end{equation}
Здесь $D_\mu Q=\partial_\mu Q+[A_{Q,\mu},Q]$, а
$A_Q=[Q,dQ]/\Delta^2$. После выбора сигнатуры относительные знаки временных
и пространственных производных фиксируются лоренцевой свёрткой.

Однако коэффициенты $Z_Q,Z_F$ и общий размерный множитель не следуют из
самой формулы составной связности. Более того, $A_Q$ остаётся связностью
косета $SO(3)/O(2)$, а не независимо выведенным полным калибровочным полем.
Поэтому \eqref{eq:v6-bosonic-lorentz-admissible-action} является допустимым
ковариантным продолжением доказанной статики, но пока не единственным
родительским действием.

\begin{proposition}[Ковариантная форма без абсолютной нормировки]
Статические члены Тома VI допускают лоренц-ковариантное объединение, но
текущая конечная архитектура не фиксирует временную кинетическую норму и
отношение двух-, четырёхпроизводного и потенциального секторов.
\end{proposition}

\section{Двухмасштабная вырожденность}

Для гладкого профиля
\begin{equation}
 q(r)=\Delta\frac{r^2}{r^2+R^2}
\end{equation}
предыдущий аудит вычислил при единичном радиусе безразмерные интегралы
\begin{align}
 I_D&=12.4786366011\ldots,\\
 I_F&=33.5371715041\ldots,\\
 I_V&=17.5002220831\ldots.
\end{align}
Общая статическая энергия имеет вид
\begin{equation}
 E(R)=c_DI_DR+c_F\frac{I_F}{R}+c_VI_VR^3.
 \label{eq:v6-bosonic-parent-scale-energy}
\end{equation}

Введём произвольные единицы длины $L_0$ и энергии $E_0$:
\begin{equation}
 c_D=\frac{E_0}{L_0},\qquad
 c_F=E_0L_0,\qquad
 c_V=\frac{E_0}{L_0^3},\qquad R=L_0r.
\end{equation}
Тогда
\begin{equation}
 \frac{E}{E_0}=I_Dr+\frac{I_F}{r}+I_Vr^3
 \label{eq:v6-bosonic-dimensionless-energy}
\end{equation}
не зависит ни от $L_0$, ни от $E_0$. Безразмерный минимум равен
\begin{equation}
 r_*=0.8301754154\ldots,
 \qquad
 \frac{M_*}{E_0}=60.7698954459\ldots.
\end{equation}
Но физические величины остаются
\begin{equation}
 R_*=L_0r_*,\qquad M_*=E_0\,60.7698954459\ldots.
\end{equation}

\begin{theorem}[Два отсутствующих размерных входа]
Текущий родитель фиксирует форму пробного профиля и его безразмерный
стационарный радиус, но не фиксирует единицы длины $L_0$ и энергии $E_0$.
Поэтому абсолютные радиус и масса не могут быть выведены из одной щели
$\Delta$, топологического класса $+15/-15$ и нормированных конечных следов.
\end{theorem}

Это не численная неопределённость. Три машинных примера с различными
$(L_0,E_0)$ воспроизводят один безразмерный профиль при произвольно разных
физических радиусах и массах.

\section{Что именно связывает функционал M300 с Хиггсом}

Единый функционал Ходжа Тома V содержит средний блок
\begin{equation}
 \Tr\bigl(XX^\dagger-|H|^2I_3\bigr)^2.
\end{equation}
При единственном уже построенном отождествлении проекторного состояния с
нормированной матрицей Грама разложим её на общую амплитуду и состояние:
\begin{equation}
 XX^\dagger=TR,
 \qquad \Tr R=1,
 \qquad Q=R-\frac{I_3}{3}.
\end{equation}
Тогда блок равен
\begin{equation}
 T^2\Tr(R^2)-2T|H|^2+3|H|^4.
 \label{eq:v6-m300-radial-higgs-only}
\end{equation}
Смешанный член зависит от радиальной амплитуды $T$, но не от
$\Tr(Q^2)$. Машинная смешанная разность по форме $Q$ и $|H|^2$ равна
$1.1\cdot10^{-16}$, тогда как смешанная производная по $T$ и $|H|^2$
равна $-2$ с остатком менее $10^{-5}$.

Если же считать позднее поле $Q$ независимым от матрицы Грама $X$, оно
вообще не входит в функционал $M_{300}$, и тот же коэффициент также равен
нулю. Обе допустимые интерпретации минимального родителя дают одинаковый
ответ.

\begin{theorem}[Нулевой портал формы в минимальном родителе]
В едином функционале $M_{300}$ коэффициент оператора
\begin{equation}
 \Tr(Q^2)H^\dagger H
\end{equation}
равен нулю. Функционал связывает поле Хиггса только с общей амплитудой
моста $T$, а не с бесследовым проекторным параметром порядка $Q$.
\end{theorem}

Следовательно, ранее найденный квадратичный портал остаётся разрешённым
эффективным оператором, но не является скрытым членом уже принятого
минимального действия.

\begin{nogocriterion}[Запрет феноменологического масштаба и портала]
Нельзя выбирать $L_0$, $E_0$ или ненулевой коэффициент
$\Tr(Q^2)H^\dagger H$ по желаемой массе, размеру, времени жизни или
космологической плотности. Переход к наблюдаемой физике допустим только
после появления нового родительского размерного оператора либо доказанного
механизма, связывающего $Q$ с радиальной амплитудой $T$.
\end{nogocriterion}

\section{Вердикт и следующая проверка}

Отрицательный результат существенно сужает программу. Бозонный дефект не
исчезает: его топология, конечный пробный профиль, составная связность и
спин ноль сохраняются. Закрывается более сильное утверждение, что текущий
родитель уже предсказывает его абсолютную массу или ненулевое взаимодействие
с Хиггсом.

Следующая проверка должна вернуться к полным уравнениям поля при общих
положительных безразмерных коэффициентах. Она обязана установить
существование стационарного решения, спектр его малых возмущений и каналы
распада независимо от абсолютного выбора $L_0,E_0$. Это отделит настоящую
безразмерную устойчивость от артефакта пробного профиля.

\begin{protocol}
\begin{enumerate}
 \item лоренц-ковариантная форма действия: \textbf{допустима};
 \item единая временная нормировка из родителя: \textbf{не выведена};
 \item безразмерные интегралы профиля: \textbf{вычислены};
 \item безразмерный стационарный радиус: \textbf{$0.8301754154\ldots$};
 \item абсолютная единица длины: \textbf{не выведена};
 \item абсолютная единица энергии: \textbf{не выведена};
 \item абсолютные масса и радиус: \textbf{не получены};
 \item радиальная связь $T|H|^2$ в $M_{300}$: \textbf{есть};
 \item портал $\Tr(Q^2)H^\dagger H$ в $M_{300}$: \textbf{коэффициент ноль};
 \item наблюдаемая связь бозонного дефекта: \textbf{не выведена};
 \item следующая проверка:
 \textbf{полные уравнения Эйлера--Лагранжа и устойчивость};
 \item рождение наблюдаемой материи: \textbf{не закрыто}.
\end{enumerate}
\end{protocol}

Машинный сертификат сохранён в
\path{s2t_v6_bosonic_defect_mass_portal_parent_gate_results.json}.